% TOP_PROBLEM: {"id":30007629,"problem_number":"LOCAL-30007629","title":"Cross-border data governance","category_id":21,"category":{"id":21,"name":"economics","display_name":"Economics","description":"Open economic theory statements, published research agendas, and proposed scoped specifications; question provenance and historical priority are qualified per record.","slug":"economics","order_index":21,"created_at":"2026-10-08T00:00:00Z"},"status":"provenance_pending","external_url":null,"legacy_ids":[],"published":false,"statement":"In the explicitly specified setting: What data-flow rules best balance digital-trade gains against privacy, cybersecurity and trust? The mathematical statement above fixes the target, assumptions and scope of this catalogue formulation.","economics":{"original_id":118,"field":"Trade, globalization and geoeconomics","kind":"design","formalization_basis":"proposed_specification","source_status":"not_verified","priority_status":"not_established","priority_note":"No primary passage stating this exact catalogue question as open was verified; no original-priority claim is made.","earliest_verified_reference":null,"current_reference":{"authors":["OECD","World Trade Organization"],"title":"Economic Implications of Data Regulation: Balancing Openness and Trust","year":2025,"url":"https://www.oecd.org/en/publications/economic-implications-of-data-regulation_aa285504-en.html","doi":"10.1787/aa285504-en","locator":"Report landing-page abstract and citation metadata","evidence_summary":"Evaluates openness and trust scenarios; no explicit optimal-governance open statement verified."},"formalization_note":"The WTO/OECD report evaluates data-policy scenarios, but an exact open optimal-governance statement was not verified. The welfare objective and privacy valuations are proposed."},"statusQualification":"Provenance pending: an explicit published statement of this exact open question was not verified. This is a catalogue-authored candidate specification, not a certified open conjecture. The WTO/OECD report evaluates data-policy scenarios, but an exact open optimal-governance statement was not verified. The welfare objective and privacy valuations are proposed.","statusReviewedAt":"2026-10-08"}
% FIRST_POSTED: 2026-10-08
% ENTRY_KIND: statement_only
% !TeX program = lualatex
\documentclass[11pt]{article}
\usepackage{fontspec}
\usepackage[a4paper,margin=25mm]{geometry}
\usepackage{amsmath,amssymb,mathtools}
\usepackage{xurl}
\usepackage[hidelinks]{hyperref}
\newcommand{\sref}[2]{\hyperlink{source:#1:#2}{[S#2]}}
\setlength{\emergencystretch}{3em}
\begin{document}
% problemId:30007629
\section*{Cross-border data governance}\label{30007629}
\subsection{Definitions and mathematical statement}
Fix jurisdictions trading digitally delivered services and transmitting personal or business data. A governance regime \(a\) specifies transfer eligibility, localization, privacy safeguards, cybersecurity requirements and enforcement. Firms and users choose participation, security effort and service use in response. Let \(T(a)\) be trade-related surplus, \(P(a)\) expected privacy loss, \(S(a)\) cyber incident loss and \(C(a)\) compliance resource cost, each defined for a declared population and under a stated probability law. The task is to characterize feasible regimes maximizing
\[W(a)=T(a)-P(a)-S(a)-C(a),\qquad a\in A.\]
The set \(A\) includes sovereignty, legal and enforceability constraints, and valuations of privacy losses must be explicitly specified rather than inferred from trade volumes. Identify how safeguards affect trust and participation separately from direct changes in transaction costs. Model cross-border externalities and regulatory interaction; unilateral optimization need not yield the jointly best regime. Estimate or bound responses using policy variation with credible controls for concurrent digital development. Report welfare sensitivity to risk valuations and heterogeneous users. Complete liberalization or localization cannot be declared optimal merely by omitting trust, privacy or cybersecurity channels from the objective.

\par\textbf{Formulation and scope.} The WTO/OECD report evaluates data-policy scenarios, but an exact open optimal-governance statement was not verified. The welfare objective and privacy valuations are proposed.

\par\textbf{Open-question provenance.} An explicit published open-question statement was not verified. Sources below are supporting literature and must not be treated as a first listing.

\par\textbf{Historical priority.} No primary passage stating this exact catalogue question as open was verified; no original-priority claim is made.

\subsection{Short English statement}
In the explicitly specified setting: What data-flow rules best balance digital-trade gains against privacy, cybersecurity and trust? The mathematical statement above fixes the target, assumptions and scope of this catalogue formulation.

\subsection{Sources}
\begin{itemize}\raggedright
\item[\textnormal{[S1]}]\hypertarget{source:30007629:1}{} OECD, World Trade Organization. 2025. Economic Implications of Data Regulation: Balancing Openness and Trust. Report landing-page abstract and citation metadata.
\url{https://www.oecd.org/en/publications/economic-implications-of-data-regulation_aa285504-en.html}.
DOI: 10.1787/aa285504-en.
{\small\textit{Supporting or current literature; see evidence qualification. Evaluates openness and trust scenarios; no explicit optimal-governance open statement verified.}}
\end{itemize}
\end{document}
